%% Chapter 2 - Literature Review and Related Work
%% Self-contained chapter file. All sections (2.1 Review method and scope;
%% 2.2 Domain and theoretical foundations; 2.3 Related research;
%% 2.4 Existing systems, products and techniques; 2.5 Gap synthesis and
%% conceptual framework) are inlined below. No external section files are
%% required. Section labels: sec:domain-foundations (2.2),
%% sec:related-research (2.3), sec:existing-systems (2.4),
%% sec:gap-synthesis (2.5).
\chapter{Literature Review and Related Work}\label{ch:literature-review}

\section{Review Method and Scope}\label{sec:review-method-and-scope}

This chapter reviews the literature needed to support the research questions
posed in Section~\ref{sec:research-questions}. The review is deliberately
targeted rather than bibliometric: its purpose is to establish (1) how
language-model copilots for security operations behave and how they validate
their own output; (2) the attack surface of applications that integrate
language models with tooling; (3) the available methods for detecting model
fabrication and for measuring answer agreement or uncertainty; and (4) which of
these techniques remain viable under the deployment constraints adopted in this
dissertation---an offline, CPU-only, single-inference budget with no retrieval
at query time. These four questions map directly onto RQ1--RQ3.

The corpus is drawn from the computing literature and from authoritative
technical sources. Peer-reviewed and established-venue material was sought
across the ACM proceedings (including the ESEC/FSE joint conference and the
Artificial Intelligence and Security (AISec) workshop), the empirical NLP and
machine learning conferences (EMNLP, ICLR, NeurIPS), the ACM Computing Surveys
journal, and the arXiv preprint archive; standards and tooling references were taken
from their canonical sources (the IETF RFC series, the OWASP project, and the
official documentation of the Nmap and Metasploit tools). Candidate works were
located with controlled query themes covering the joint areas of interest:
security copilots and penetration testing assisted by language models, prompt
injection and indirect prompt injection, hallucination detection, semantic
uncertainty, self-consistent decoding, retrieval-augmented generation, and
authorisation control for tool-driven systems.

Works were included when they satisfied two criteria: relevance to at least one
research question, and a verifiable primary source reporting either an empirical
evaluation or an authoritative specification. The machine-learning literature
covered is concentrated in the period 2019--2025, which spans the advent of the
sentence-embedding and language-model techniques relied on here
\parencite{reimers2019sbert,lewis2020rag}; older technical standards and
tooling references are included regardless of year because they define stable
interfaces rather than evolving research claims. Preprint versions were
preferred only when a peer-reviewed version was not available for a work; each
entry in the reference list was verified against its real source before it was
cited in this dissertation. The review is bounded by the deployment constraint:
methods that assume interactive generation, a cloud-hosted verifier, or
retrieval infrastructure are surveyed because they are the competing
alternatives against which the guarded pipeline must be justified, and their
unsuitability for the offline single-inference budget is part of the analysis
in Section~\ref{sec:gap-synthesis}.

The remainder of the chapter is organised as follows.
Section~\ref{sec:domain-foundations} establishes the theoretical background of
intent recognition and of the guarantees expected from a safety-critical
interface. Section~\ref{sec:related-research} examines the related research on
LLM copilots, attacks, and fabrication detection in detail.
Section~\ref{sec:existing-systems} reviews existing systems, products and
techniques against the same criteria. Section~\ref{sec:gap-synthesis} synthesises
the findings into the gap that motivates this dissertation and into the
conceptual framework that the design chapters build on.

\TODO{Confirm whether the faculty's assessment rubric requires a formal
systematic-review protocol (screening counts, PRISMA-style reporting). The
current review is a targeted narrative review; screening statistics are not
reported because no formal screening process was conducted.}

\section{Domain and Theoretical Foundations}\label{sec:domain-foundations}

\subsection*{The security-operations domain}

The domain this dissertation addresses is offensive and defensive security
operations: the systematic discovery, assessment and (where authorised)
remediation of weaknesses in networks and applications. The operator's work is
organised around a small set of recurring intents---discovering what is
exposed, enumerating services, identifying known vulnerabilities, and checking
whether specific weaknesses apply to particular targets. These intents map onto
specialised tooling: network discovery and scanning (of which Nmap is the
canonical example \parencite{lyon2009nmap}), and exploitation framework
drivers (of which Metasploit is the canonical example \parencite{metasploit}).
Two properties of this domain matter for the work in this dissertation. First,
each intent carries a \emph{target} and a set of \emph{parameters}, so the gap
between prose and machine dispatch is structural rather than cosmetic. Second,
every action is bounded by authorisation: a scan may only touch targets and
techniques within an explicitly declared engagement scope, and some intents are
reserved for roles with the required mandate. The translation layer evaluated
in this dissertation is exactly the component that must preserve both the
structure and the authorisation of operator intent.

\subsection*{Intent recognition and natural-language understanding}

At its core, intent recognition is a natural-language understanding (NLU)
problem: given an utterance, determine the communicative intent and extract the
arguments that parameterise it. In a classical NLU system this is split into an
intent-classification step, which assigns the utterance to one of a fixed set
of intent classes, and a slot-filling (or entity-extraction) step, which
recovers typed argument values such as hosts, networks and ports. The intent
contract introduced in Section~\ref{sec:individual-component-context} is the
machine-consumable record of these two steps: a resolved intent type, a target
and recognised options, together with any validation findings. For this record
to be useful to downstream automation it must be expressed in a
schema-bounded, machine-parseable syntax, which is why the IRE emits a
structured contract rather than prose.

Representational infrastructure underpins the more advanced parts of the
design. Sentence embeddings map natural-language strings into a vector space in
which semantically similar strings are close under a distance such as cosine
similarity; Sentence-BERT is the representative technique \parencite{reimers2019sbert}.
These embeddings provide the two similarity signals used in this work: the
agreement signal by which multiple model outputs are judged to express the same
meaning, and the cache signal by which a repeated intent is recognised without
re-invoking the model.

\subsection*{Language models as stochastic components}

Instruction-tuned language models are the generative core of the IRE. Their
defining property for this dissertation is that inference is a sampling
process: given a prompt, the model defines a probability distribution over
output sequences, and a response is drawn from that distribution. Sampling
temperature controls how peaked the distribution is, so the same prompt can
yield different outputs across calls. A single response is therefore a
\emph{sample} of the model's belief, not a deterministic verdict. This
observation is the theoretical basis of the design's treatment of the model as
an unreliable proposition generator whose output must be checked rather than
trusted. It also underlies consistency-based confidence: the
self-consistency decoding strategy samples several responses and strengthens
confidence in answers that agree with one another \parencite{wang2023selfconsistency},
applying a voting intuition that originates in chain-of-thought reasoning and
transfers naturally to classification agreement.

\subsection*{Hallucination as a failure class}

A second property of language models is that their outputs can be fabricated
relative to the true state of the world. The hallucinations relevant here are
the \emph{extrinsic} form in the taxonomy of Ji \emph{et al.}: statements that
contradict or are unverifiable against an authoritative source, rather than
mere inconsistency with an input context \parencite{ji2023hallucination}.
Fabricated security identifiers are a concrete instance: a model may emit a
CVE identifier, an IP address or a port that does not exist or is not the one
the operator intended. Because a downstream tool will act on those values, the
practical requirement is not simply to detect fabrication in general but to
detect it in tightly constrained value spaces (identifier formats, address
syntaxes, port ranges), where deterministic checking is feasible.

\subsection*{Agreement and uncertainty as fabrication signals}

Detection without an external verifier relies on the statistical structure of
the model's own output. Sampling-based methods compute a consistency score from
a small set of generated responses---SelfCheckGPT scores sampled responses for
mutual consistency \parencite{manakul2023selfcheckgpt}---while information-theoretic
approaches frame the uncertainty of free-form answers through semantic
entropy over meaning-equivalence clusters \parencite{kuhn2023semantic,
farquhar2024semantic}. These methods provide the theoretical grounding for the
self-consistency gate in the IRE design. Their common cost, however, is
repeated generation: every extra sample multiplies inference time, which
collides with the CPU-only, single-user latency budget adopted in this
dissertation. That tension---between the statistical value of sampling and the
hard cost of each sample---is a central theme returned to in the synthesis
(Section~\ref{sec:gap-synthesis}).

\subsection*{Authorisation and network boundaries}

Finally, two security foundations bound what the validated contract is allowed
to assert. The first is role-based access control (RBAC), under which a user's
authorisation to perform an action is determined by the roles assigned to that
user and an access matrix that grants or denies the intent for each role; the
IRE's four-role model operationalises this standard pattern. The second is
network-level scope: private address allocation is codified by RFC~1918
\parencite{rfc1918}, which defines the address ranges reserved for private
internets and provides the technical basis for distinguishing an internal
target from an external one. Together with the operator-declared engagement
scope, these foundations define the predicates that deterministic validation
can and cannot decide: syntax, format and authorisation are decidable
propositions; only meaning is delegated to the model. This division of
responsibility is the theoretical spine of the solution designed in
Chapter~\ref{ch:product-design}.

\TODO{If the assessment rubric requires canonical textbook references for NLU
fundamentals and for role-based access control, add formal citations for the
intent-classification/slot-filling formulation and for the RBAC reference model;
none are currently in the reference database.}

\section{Related Research}\label{sec:related-research}

This section reviews the three bodies of work most directly related to the
Intent Recognition Engine---language-model copilots for security operations,
attacks on applications that integrate language models, and methods for
detecting model fabrication---followed by a brief assessment of what each body
contributes to and leaves open for this dissertation.

\subsection*{Language-model copilots for security operations}

Two systems define the current shape of LLM-assisted offensive security. Happe
and Cito proposed language models as \emph{sparring partners} for penetration
testing and demonstrated a closed-loop system in which a model drove a
deliberately vulnerable virtual machine to root access \parencite{happe2023pwned}.
PENTESTGPT treats penetration testing as a task for a language-model-centric
assistant, decomposing an engagement across a reasoning chain, a generation
chain and a guidance chain so that each phase of an assessment is planned and
executed incrementally \parencite{deng2024pentestgpt}. Both contributions
establish that language models can produce \emph{plausible} security actions
and can operate tooling in controlled environments, and both are representative
of the generation-side focus of the field: output is produced for a human
operator to review, and correctness is handled informally. Neither system
subjects its output to deterministic schema, scope and authorisation validation
before machine consumption, which is precisely the gap that a copilot feeding
an \emph{automated} toolchain must close. The IRE inverts this priority: it
treats the model's generation as a candidate that must pass checkable
predicates (Section~\ref{sec:domain-foundations}) and only then becomes an
actionable intent.

\subsection*{Attacks on language-model-integrated applications}

The second body of work establishes that unvalidated model output is dangerous
in integrated systems. Liu \emph{et al.} formalised prompt-injection attacks
against LLM-integrated applications and systematically benchmarked a set of
common defences, finding that prevention-based defences were weak and
detection-based defences incomplete \parencite{liu2023promptinjection}.
Greshake \emph{et al.} demonstrated \emph{indirect} prompt injection, in which
an attacker plants instructions inside content---including web pages and
documents---that an integration later retrieves and feeds to the model, and
showed compromise of real-world assistant-style integrations
\parencite{greshake2023indirect}. Kang \emph{et al.} showed that dual-use
attacks that combine well-known computer-security techniques bypass in-the-wild
content safeguards of deployed models with near-complete success
\parencite{kang2024dualuse}. The OWASP Top~10 for LLM applications
consolidates these classes---prompt injection first among them---as
first-class application risks rather than theoretical concerns
\parencite{owasp2025}.

The implication for this dissertation is direct. A security copilot that
consumes an utterance, asks a model to classify it, and dispatches the result
is exactly the kind of integration these attacks exploit: a single adversarial
utterance can steer the classification output, and any downstream consumer will
then act on the attacker-chosen intent unless an independent, non-model check
intervenes. The relevant design response therefore places deterministic
authorisation and validation \emph{after} generation, so that even a perfectly
diverted model cannot manufacture an accepted intent that violates syntax,
scope or role policy (the safety invariant introduced in
Chapter~\ref{ch:product-design}).

\subsection*{Detecting fabrication and measuring uncertainty}

The third body of work supplies the theoretical toolkit for catching fabricated
output. Ji \emph{et al.} provide the taxonomy of hallucination in natural
language generation that distinguishes intrinsic from extrinsic fabrication
\parencite{ji2023hallucination} (Section~\ref{sec:domain-foundations}).
Sampling-based detection uses the model's own output statistics:
SelfCheckGPT scores the consistency of sampled responses to flag ungrounded
claims \parencite{manakul2023selfcheckgpt}; semantic-uncertainty methods embed
responses and cluster meaning-equivalent outputs to measure disagreement
\parencite{kuhn2023semantic}, and semantic entropy applies this principle to
detect confabulations in free-form text \parencite{farquhar2024semantic}.
Self-consistency decoding accumulates agreement across samples and rewards the
majority answer \parencite{wang2023selfconsistency}.

What these methods share is an implicit budget: they assume either interactive
generation with a user, or the freedom to produce several responses per query
and cluster them. Each additional sample multiplies the inference cost of a
CPU-only model on the operator's laptop, so the design of this dissertation
must transpose the disagreement intuition into a \emph{single} guarded
classification: restricted sampling (a bounded number of responses), a cheap
embedding-based agreement gate, and a deterministic escalation path that does
not depend on a large verifier model.

\subsection*{Grounding via retrieval and embeddings}

A complementary line of work grounds model answers in external corpora rather
than in self-consistency. Retrieval-augmented generation (RAG) conditions
generation on retrieved passages so that answers can be traced to sources
\parencite{lewis2020rag}, and sentence embeddings provide the similarity
signals such systems rely on \parencite{reimers2019sbert}. RAG is attractive
for factual Q\&A, but it introduces retrieved content into the prompt at
inference time---an additional prompt-injection surface
\parencite{greshake2023indirect}---and presupposes corpus infrastructure and
embedding services that contradict the offline, single-device boundary adopted
here (Section~\ref{sec:scope-limitations-assumptions}). The IRE therefore
performs no retrieval at inference time and uses embeddings only for the
agreement and cache signals.

\subsection*{Assessment}

Taken together, the related research establishes three things: copilots can
generate plausible security actions but do not validate them for machine
consumption; integrated systems are exploitable through adversarial and
indirect-injection utterances, and published defences are brittle; and
fabrication can be detected in principle through sampling and semantic
agreement, at a cost that grows with the number of samples. What the surveyed
literature does not measure, to the best of the review conducted here, is the
behaviour of a \emph{guardrail stack}---a full deterministic validation
pipeline with authorisation checks and a bounded self-consistency gate over a
small offline model---under a strict single-inference latency budget, nor the
isolated contribution of each guard layer, nor whether the achieved safety
transfers between different offline model backends. Those measurements are
precisely RQ1--RQ3, and Section~\ref{sec:gap-synthesis} consolidates the gap
they define into the conceptual framework used in the remainder of this
dissertation.

\TODO{At the final AUDIT step, re-verify the characterisation of
\citeauthor{liu2023promptinjection} (defence benchmark) and
\citeauthor{kang2024dualuse} (dual-use bypass) against the full text of their
respective publications, since the descriptions here are relayed through the
project's research paper rather than a direct reading of each source.}

\section{Existing Systems, Products and Techniques}\label{sec:existing-systems}

This section examines the systems, products and techniques that currently occupy
the space in which the Intent Recognition Engine operates. They are reviewed
against the same inclusion criteria as the wider survey (Section~\ref{sec:review-method-and-scope}): relevance to at
least one research question and a verifiable primary source reporting either an
empirical evaluation or an authoritative specification. Where a surveyed option
is a technique rather than a delivered product, the assessment asks whether it
can be combined with the deployment constraint adopted in this
dissertation---offline, CPU-only, single-inference, retrieval-free operation.

\subsection*{Deterministic security tooling}

The techniques that the IRE ultimately drives are mature and deterministic. Nmap
is the canonical network-discovery and scanning tool, whose techniques---host
discovery, port enumeration, service and version detection---define the target
values an intent contract must carry \parencite{lyon2009nmap}. Metasploit is the
canonical exploitation-framework driver, whose modules and options parameterise
the most privileged class of intents \parencite{metasploit}. Both expose
structured command interfaces (a tool name, target, options and constraints)
rather than free-form prose, and neither embeds a natural-language layer of its
own. In current practice the mapping from operator intent onto these structured
interfaces is performed by a human analyst or by a scripted wrapper, and any
authorisation policy has to be wired externally; RFC~1918 \parencite{rfc1918}
provides only the technical basis for distinguishing internal from external
addressing. These tools therefore constitute the deterministic ``ground truth''
that a translation layer must produce---and the reason a translation error
(Chapter~\ref{ch:product-design}) becomes a machine-dispatched error.

\subsection*{Language-model copilots for security}

The recent systems that add natural language to this toolchain are generation-side
copilots. PENTESTGPT decomposes a penetration-testing engagement across reasoning,
generation and guidance chains and executes each phase incrementally, producing
structured guidance that a human operator reviews \parencite{deng2024pentestgpt}.
The sparring-partner system demonstrated by Happe and Cito is a closed loop in
which a model proposes low-level actions that are executed automatically against
a deliberately vulnerable machine, up to root access \parencite{happe2023pwned}.
Both establish that a model can drive security tooling in controlled
environments, and both treat output correctness informally: the intended consumer
of the output is human, and neither system subjects generation to deterministic
schema, scope or authorisation validation before an action is taken. As products,
they are representative of the ``AI assistant'' category of security copilots
rather than of an automated, machine-consumed intent layer.

\subsection*{Model-integration techniques: retrieval and tool calling}

Beyond security copilots, two integration techniques are common to LLM-driven
automation generally. Retrieval-augmented generation (RAG) conditions answers on
external corpora so that claims can be traced to sources
\parencite{lewis2020rag}, using sentence embeddings as the similarity signal
\parencite{reimers2019sbert}; tool-calling scaffolds connect a model to external
commands through function-call interfaces. As deployed techniques both widen the
integration surface: retrieved content enters the prompt at inference time and is
therefore an indirect-injection vector \parencite{greshake2023indirect}, and
tool-calling binds a model action to a machine effect with no intermediate
validation stage. The IRE is deliberately retrieval-free and does not render
free-form model output to a downstream shell, which narrows these channels to the
classification of a single utterance (Section~\ref{sec:domain-foundations}).

\subsection*{Guardrail techniques for model-integrated applications}

At the application layer, guardrail practice is codified by the OWASP Top~10 for
LLM applications, which catalogues prompt injection (LLM01), insecure output
handling and related classes as first-class risks rather than theoretical
concerns \parencite{owasp2025}. The defence techniques surveyed in
Section~\ref{sec:related-research} are the prevention-detection families
benchmarked by Liu \emph{et al.} \parencite{liu2023promptinjection}---shown in
that study to be individually insufficient---together with detection techniques
for indirect injection \parencite{greshake2023indirect} and the empirically
demonstrated bypass of deployed content safeguards by dual-use attacks
\parencite{kang2024dualuse}. What these techniques have in common is that they
act as app-layer checks over a model's behaviour after the fact. None of them is
integrated as a deterministic validation stack that sits \emph{between} model
generation and machine dispatch and that makes syntax, scope and authorisation
independently decidable (Section~\ref{sec:domain-foundations}).

\subsection*{Fabrication-detection and uncertainty techniques}

A separate product-facing technique family detects fabricated output. SelfCheckGPT
scores the mutual consistency of sampled responses \parencite{manakul2023selfcheckgpt};
semantic-uncertainty and semantic-entropy methods measure answer disagreement over
meaning-equivalence clusters \parencite{kuhn2023semantic,farquhar2024semantic};
and self-consistency decoding accumulates agreement across samples
\parencite{wang2023selfconsistency}. These are deployed as techniques that assume
repeated or interactive generation. Two external references complement them for
the value spaces relevant here: the National Vulnerability Database (NVD) is the
authoritative registry of CVE identifiers \parencite{nistnvd}, and the Exploit
Prediction Scoring System (EPSS) scores the likelihood that a vulnerability will
be exploited \parencite{firstepss}. Both are the deterministic, offline-unfriendly
references that a fabricated-identifier guard would consult if the offline
boundary were relaxed (Chapter~\ref{ch:results}).

\subsection*{Comparative assessment}

Table~\ref{tab:existing-systems} consolidates the judgments from the subsections above against the properties that matter for the RQ1--RQ3
measurement (Section~\ref{sec:research-questions}). ``Schema check'' and
``scope/authorisation enforcement'' record whether a system subjects
machine-consumed output to structural and policy checks; ``offline CPU'' and
``retrieval-free'' record whether the system can run under the deployment
constraint adopted here.

\begin{table}[h]
\centering
\footnotesize
\caption{Surveyed existing systems and techniques against the properties adopted in this dissertation. Judgements summarise the cited primary sources and are re-verified at the final AUDIT.}
\label{tab:existing-systems}
\begin{tabular}{@{}lcccc@{}}
\hline
\textbf{System / technique} & \textbf{Human review} & \textbf{Schema check} & \textbf{Scope / authz} & \shortstack{\textbf{Offline CPU}\\\textbf{retrieval-free}} \\
\hline
Nmap / Metasploit & -- & Yes (interface) & External only & Yes \\
PENTESTGPT & Yes & No & No & -- \\
Sparring partner & Partial & No & No & -- \\
RAG agents & Partial & No & No & No \\
Fabrication detection & -- & No & No & Partial \\
IRE (this work) & No & Yes & Yes & Yes \\
\hline
\end{tabular}
\end{table}

In summary, the surveyed landscape contains tools that are deterministic but
require manual intent mapping, copilots that map intent but consume human review,
integration techniques that connect models to tools without an intermediate
validation stage, and defence techniques that check model behaviour after the
fact. None combines a deterministic schema, scope and authorisation validation
stack with offline, CPU-only, single-inference and retrieval-free operation over
a closed intent taxonomy. That combination is the space the existing-systems
survey leaves open and is consolidated in Section~\ref{sec:gap-synthesis}.

\TODO{If the assessment rubric requires a survey of \emph{named} commercial
security copilots (for example cloud SIEM/EDR copilot products), provide
verifiable primary sources for each; none are currently in the reference
database, so no commercial product is named here.}

\TODO{At the final AUDIT step, re-verify the comparative-table judgements in
Table~\ref{tab:existing-systems} against the full text of each cited source.}

\section{Gap Synthesis and Conceptual Framework}\label{sec:gap-synthesis}

The review in this chapter converges on three findings. First, the
language-model copilots surveyed \parencite{deng2024pentestgpt,happe2023pwned}
generate plausible security actions but do not validate them for machine
consumption: their output is produced for a human reviewer, and correctness is
handled informally (Section~\ref{sec:related-research}). Second, applications
that integrate a model with tooling are exploitable through adversarial and
indirect-injection utterances, the surveyed defences are individually
insufficient \parencite{liu2023promptinjection,greshake2023indirect}, and the
existing-systems survey shows that current guardrail techniques act as
post-hoc application-layer checks rather than as a deterministic validation
stage between generation and dispatch (Section~\ref{sec:existing-systems}).
Third, fabrication-detection and agreement methods
\parencite{manakul2023selfcheckgpt,kuhn2023semantic,farquhar2024semantic,wang2023selfconsistency}
are available in principle, but they assume sampling or interactive generation
whose cost collides with the offline, CPU-only, single-inference budget adopted
here (Section~\ref{sec:domain-foundations}).

The gap that follows from these three findings is the combination that the
existing-systems survey leaves open: no surveyed system delivers a dedicated,
isolated natural-language understanding layer that converts free-form
offensive-security instructions into validated, schema-constrained, authorised
machine contracts while running offline, retrieval-free, on a single inference
budget. This dissertation addresses exactly that combination by designing and
\emph{measuring} a guarded intent-recognition engine.

\subsection*{Conceptual framework}

The design chapters build on three principles established in
Section~\ref{sec:domain-foundations}.

\emph{First}, the language model is treated as a \emph{stochastic proposition
generator}: a single generated response is a sample weighted by the model's
beliefs rather than a deterministic verdict
(Section~\ref{sec:domain-foundations}). Nothing the model emits is trusted until
it satisfies checkable predicates; sampling-based agreement
\parencite{wang2023selfconsistency} provides a confidence signal, but the
generation itself is always a candidate rather than a decision.

\emph{Second}, responsibility is divided along the decidability boundary
identified in Section~\ref{sec:domain-foundations}: syntax, format and
authorisation are decidable predicates (schema, host and CIDR syntax, port
ranges, network and scope policy, role permission), while only meaning is
delegated to the model. The guarded pipeline therefore couples a small
instruction-tuned model with deterministic stages---schema validation, regex and
network validation, scope enforcement, role-based access control---whose purpose
is to contain model error before it can propagate to a planner or executor
(Chapter~\ref{ch:product-design}).

\emph{Third}, the design enforces a \emph{safety invariant}: an intent contract
is accepted only when it satisfies the conjunction of the deterministic security
predicates---schema-valid, network-authorized, in-scope, role-permitted and
non-adversarial. The model may propose any candidate, but it cannot force
acceptance of a candidate that violates any predicate in the conjunction. This
converts the open problem of making a single language model trustworthy into the
closed problem of classifying a single utterance under a known policy---the
reduction that justifies acting on model-produced security instructions in an
automated toolchain.

On top of these principles the framework adopts the bounded-sampling resolution
of the cost tension from Section~\ref{sec:domain-foundations}: a small, bounded
number of responses, a cheap embedding-based agreement gate, and a deterministic
escalation path (the hallucination validator) that does not depend on a large
verifier model (Section~\ref{sec:related-research}).

\subsection*{From framework to measurement}

The framework yields three measurable claims, each mapped to the corresponding
research question (Section~\ref{sec:research-questions}):
\begin{itemize}
\item \textbf{Pipeline value (RQ1).} The guarded pipeline raises the
probability that model output satisfies the full conjunction of security
predicates relative to raw generation; this is quantified by comparing intent
accuracy, flaw-catch rate and false-positive rate with and without the guards.
\item \textbf{Layer contribution (RQ2).} The contribution of each guard layer is
quantified by ablating layers in turn against a constant baseline, isolating the
largest safety contributor and the layers that trade safety for usability.
\item \textbf{Generality (RQ3).} The identical guardrail stack is applied to two
distinct offline backends to determine how much of the measured safety is
backend-independent rather than an artefact of a single model.
\end{itemize}
The evaluation protocol that operationalizes these questions is developed in
Chapter~\ref{ch:methodology}, the implementation and automated test suite in
Chapter~\ref{ch:implementation-testing}, and the measurements themselves in
Chapter~\ref{ch:results}. Together with the preceding sections, this concludes
the review, sets out the conceptual framework that guides the design of the
Intent Recognition Engine in Chapter~\ref{ch:product-design}.

\TODO{At the final AUDIT step, cross-check the informal safety invariant stated
here against the full formal definition to be developed in
Chapter~\ref{ch:product-design}.}